\subsection{Original conormal and cyclic relation
towers}\label{conormal_cyclic_pr21:original-conormal-and-cyclic-relation-towers}

\subsubsection{Scope and source reconciliation}\label{conormal_cyclic_pr21:scope-and-source-reconciliation}

This add-only continuation starts from main at \nolinkurl{a4494fba4968db837a8af6fcb952c8688cd5d1da}. It
retains the merged Codex quotient/homotopy/boundary implementation, the merged finite-frontier
implementation, and the reviewed sum-connection note. It uses the owner's attached \textbf{Handoff:
sum-generated arithmetic control}, especially targets 1, 2, 4 and 8. The current upload is an
analytic/algebraic continuation with a separately reported 24-method finite checker; it is not a
Lean certificate. The local execution service failed before its ZIP could be opened, so neither the
ZIP manifest nor that checker has been independently rerun here.

Two actual GitHub corrections are incorporated, not merely acknowledged:

\begin{itemize}
\tightlist
\item
  Commit \nolinkurl{600ab7a527f7a626c162bf0199d093d454b86c14} makes the nonempty hypothesis explicit in
  the monic-freeness argument and calculates h=1 separately. Every positive-depth relation ideal for
  h=1 is the whole polynomial ring. The corresponding underlying quotients are zero, while the split
  lift still retains tau and supported zero. The analytic theta seed and its mass do not disappear.
  \nolinkurl{ConormalTower.top_level_zero} checks the full-ideal quotient at every level.
\item
  Commit \nolinkurl{7aec03a93b66d3d97fbf6f7042f58bb2c1019543} corrects the completed integration audit
  to 68 selected declarations and enables the existing frontier workflow on main and relevant pull
  requests. This branch inherits those files unchanged. The current 68-target audit, not the
  historical 66-target description, is rerun by the added workflow.
\end{itemize}

The original scalar G, e, tau, arithmetic projection to the infinite ring, reconstructed carriers,
quotient transports and theta normalization g=2 xi are unchanged. No claim about a global arithmetic
small-control estimate follows from the present algebraic results.

\subsubsection{1. An actual derivative between consecutive original
quotients}\label{conormal_cyclic_pr21:an-actual-derivative-between-consecutive-original-quotients}

Let R and A be commutative rings, A an R-algebra, I an ideal of A, and D an R-derivation of A. Then

\begin{verbatim}
D(I^(r+1)) is contained in I^r,  r >= 0.
\end{verbatim}

Proof: use induction on r and the additive generation of an ideal product by products. For a in
I\^{}(r+1), b in I, the Leibniz terms a D(b) and b D(a) both lie in I\^{}(r+1): the first by the
ideal property, the second by the induction hypothesis. At r=0 the target ideal is the whole ring.

\nolinkurl{SplitZeroConormalTower.lean} implements that proof as \nolinkurl{deriv_mem_pow}. It then uses
the merged \nolinkurl{SplitZero.RelationLayer.derivative} to construct

\begin{verbatim}
delta_r : A/I^(r+1) ->_R A/I^r,  [x] |-> [D x].
\end{verbatim}

It separately constructs the quotient projection p\_r between the same types and proves

\begin{verbatim}
p_r delta_(r+1) = delta_r p_(r+1).
\end{verbatim}

The derivative is not an endomorphism of A/I. These are coefficient-linear maps on the actual
original quotient types, not new quotient carriers. At a retained relation z in I,

\begin{verbatim}
delta_1([z a]) = [a D(z)].
\end{verbatim}

For the full arithmetic multiplier u, the implemented formula is

\begin{verbatim}
delta_r([u x]) = [u D(x)] + [x D(u)].
\end{verbatim}

It applies in particular to the retained Taylor unit, without assigning the unit the value one or
deleting its derivative. The definitions work for arbitrary commutative rings and do not assume
characteristic zero or squarefreeness.

\subsubsection{2. The exact cyclic pullback, before a polynomial is used to describe
it}\label{conormal_cyclic_pr21:the-exact-cyclic-pullback-before-a-polynomial-is-used-to-describe-it}

Fix S in A. At each depth define the actual coefficient submodule

\begin{verbatim}
K_r = {P in R[X] : P(S) belongs to I^r}.
\end{verbatim}

The implementation calls this \nolinkurl{cyclicLevel}; it is literally the comap of the original
relation submodule along \nolinkurl{Polynomial.aeval\ S}. Its constructed quotient map is

\begin{verbatim}
alpha_r : R[X]/K_r ->_R A/I^r,  [P] |-> [P(S)].
\end{verbatim}

\nolinkurl{evaluation_injective} proves injectivity by the two original quotient equality criteria. No
assumed injectivity of polynomial evaluation is required. This implementation is a linear map; its
usual algebra-homomorphism packaging is not claimed as a new bundled declaration.

If D(S)=1, Mathlib's formal polynomial chain rule gives

\begin{verbatim}
D(P(S)) = P'(S).
\end{verbatim}

Thus P in K\_(r+1) implies P' in K\_r. The code constructs the cyclic derivative and proves the full
square

\begin{verbatim}
R[X]/K_(r+1) --d/dX--> R[X]/K_r
      | alpha_(r+1)          | alpha_r
      v                      v
   A/I^(r+1) ------D------> A/I^r.
\end{verbatim}

The actual unit comparison is also proved:

\begin{verbatim}
delta_r([u P(S)]) = [u P'(S)] + [D(u) P(S)].
\end{verbatim}

The final term need not be in the cyclic image. The equation retains its larger target. Most
importantly, K\_r is not replaced by K\_1\^{}r. The explicit annihilator calculation below explains
the difference.

\subsubsection{3. Exact local maximum, including an attaining
monomial}\label{conormal_cyclic_pr21:exact-local-maximum-including-an-attaining-monomial}

Here the base field is C. Let k\textgreater=1, d\textgreater=1, r\textgreater=1 and write the
complete packet polynomial as

\begin{verbatim}
h(s) = product_rho (s-rho)^(m_rho).
\end{verbatim}

For a root tuple bold-rho=(rho\_1,\ldots,rho\_k), put m\_i=m\_(rho\_i) and y\_i=s\_i-rho\_i. Locally
the ideal (h(s\_i)) is (y\_i\^{}(m\_i)), because each remaining factor is a unit. In the quotient by
its r-th power the monomial basis is precisely

\begin{verbatim}
y^a such that sum_i floor(a_i/m_i) < r.
\end{verbatim}

Write a\_i=m\_i q\_i+t\_i with 0\textless=t\_i\textless m\_i. Its total degree satisfies

\begin{verbatim}
sum_i a_i <= sum_i(m_i-1) + (r-1) max_i m_i.
\end{verbatim}

The proof uses sum\_i q\_i \textless= r-1 and bounds each m\_i by its maximum. This bound is
attained: choose j with m\_j=max\_i m\_i, let every t\_i=m\_i-1, let q\_j=r-1, and let every other
q\_i=0.

\nolinkurl{SplitZeroCyclicDepth.lean} formalizes the quotient/remainder predicate, the degree bound,
that explicit exponent vector, and attainment. The floor-sum formulation is linked through Euclidean
division. This is an all-index, all-depth result, not a finite calibration. The additional algebraic
identification of these monomials with the polynomial quotient and the characteristic-zero
multinomial step below have written proofs, not an additional Lean certificate in these two modules.

Let N=y\_1+\ldots+y\_k. Its exact nilpotency index is

\begin{verbatim}
L_r(bold-rho) = 1 + sum_i(m_i-1) + (r-1) max_i m_i.
\end{verbatim}

Indeed every monomial of total degree L\_r is killed. The attaining monomial at degree L\_r-1 has
coefficient

\begin{verbatim}
(L_r-1)! / product_i a_i!
\end{verbatim}

in N\^{}(L\_r-1), a nonzero positive integer in characteristic zero. Distinct surviving monomials
cannot cancel. This proves both vanishing and minimality.

The proof also counts the local dimension as

\begin{verbatim}
(product_i m_i) * binomial(k+r-1,k).
\end{verbatim}

The q\_i with sum less than r give the binomial factor, and their remainders give the product
factor.

\subsubsection{4. Collisions require a maximum at every
depth}\label{conormal_cyclic_pr21:collisions-require-a-maximum-at-every-depth}

Set S=sum\_i s\_i. At a collided sum lambda, take the maximum over all root tuples with that sum:

\begin{verbatim}
ell_(lambda,r) = max_(sum rho_i=lambda) L_r(bold-rho).
chi_r(X) = product_lambda (X-lambda)^(ell_(lambda,r)).
\end{verbatim}

The Chinese-remainder decomposition and the exact local nilpotency prove

\begin{verbatim}
I^r intersect C[S] = (chi_r(S)).
\end{verbatim}

The same annihilator is obtained in the variable-permutation invariant subalgebra because every
polynomial in S is already invariant. Thus alpha\_r identifies the cyclic algebra C{[}X{]}/(chi\_r)
with its actual image, with no squarefree simplification.

The maximizing tuple can change with depth. In the declared calibration packet with roots 0,1,2 of
multiplicities 5,4,1, respectively, the sum lambda=2 in tensor degree two has

\begin{verbatim}
ell_(2,r) = max(5r,4r+3).
\end{verbatim}

At depths 1,\ldots,5 the orders are 7,11,15,20,25. A tuple chosen only at depth one is insufficient
for the whole tower. This is an algebraic test packet, not a statement about actual zeta zeros.

For every lambda,

\begin{verbatim}
ell_(lambda,r+1) >= ell_(lambda,r)+1,
ell_(lambda,r) <= r ell_(lambda,1).
\end{verbatim}

The first follows by retaining a maximizing tuple from depth r; its next slope is a positive
multiplicity. For the second, note max\_i m\_i \textless= 1+sum\_i(m\_i-1) for every tuple. Hence

\begin{verbatim}
chi_r divides chi_(r+1),
chi_r divides chi_(r+1)',
chi_r divides chi_1^r.
\end{verbatim}

The last divisibility supplies the stated surjection C{[}X{]}/(chi\_1\^{}r) -\textgreater{}
C{[}X{]}/(chi\_r), not an equality. For h(s)=(s-rho)\^{}2 and k=2,

\begin{verbatim}
chi_r(X)=(X-2rho)^(2r+1),
chi_1(X)^r=(X-2rho)^(3r).
\end{verbatim}

They differ as soon as r\textgreater1. The derivative maps C{[}X{]}/chi\_(r+1) -\textgreater{}
C{[}X{]}/chi\_r are onto over C: integrate any polynomial representative, then take its source
class. Their kernel dimensions are deg(chi\_(r+1))-deg(chi\_r).

For the empty packet h=1, handled separately as required by the GitHub correction, I=P,
K\_r=C{[}X{]}, and chi\_r=1 for all positive depths. The cyclic arithmetic quotient is zero. No
maximum over an empty tuple set or inverse of a zero-dimensional Gram matrix is used to define it.
The original nonzero analytic seed remains available as an analytic object.

\subsubsection{5. Explicit equivariant retraction of the cyclic summand (written
proof)}\label{conormal_cyclic_pr21:explicit-equivariant-retraction-of-the-cyclic-summand-written-proof}

This checks the handoff's retraction formula, including its inverse factor. On a lambda-primary part
M of the symmetric finite algebra, let N=S-lambda, N\^{}ell=0, and let w be the lambda-component of
the retained arithmetic unit U. Exact maximality gives N\^{}(ell-1)w != 0. Choose the stated
top-monomial coefficient functional theta with theta(N\^{}(ell-1)w) != 0. If it is realized first on
an ordered root component, it can be restricted to the invariant subspace; the component of the
symmetric vector is still detected.

Define

\begin{verbatim}
pi0(v) = sum_(j=0)^(ell-1) theta(N^(ell-1-j)v) X^j mod X^ell.
\end{verbatim}

Coefficient comparison gives pi0(Nv)=X pi0(v): the constant term on the left is theta(N\^{}ell v)=0
and the other coefficients shift exactly. Put c=pi0(w). Its constant coefficient is nonzero, so it
is invertible modulo X\^{}ell. Explicitly, if c=c0(1+n), the inverse is c0\^{}(-1)
sum\_(j=0)\^{}(ell-1)(-n)\^{}j.

Then pi=c\^{}(-1)pi0 satisfies

\begin{verbatim}
pi(P(N)w)=P(X) mod X^ell.
\end{verbatim}

Therefore the unit-weighted cyclic injection eta has an S-equivariant left inverse. Summing over
lambda constructs the full retraction, and ker(pi) is an invariant complement. This is a module
splitting, not an assertion that pi is an algebra map or that the summands are orthogonal for the
canonical arithmetic Gram. All cross terms under that splitting remain.

This cyclic summand retains each amplified value k rho; collided repeated pure tensors need not
equal the particular cyclic eigenvector. The cyclic residue trace and the full symmetric trace have
their distinct multiplicities. No part of the present calculation identifies those traces or drops
the U-dagger U factor in the ambient arithmetic pairing.

\subsubsection{6. Formal versus analytic scope}\label{conormal_cyclic_pr21:formal-versus-analytic-scope}

The new formal outputs are the ideal-power derivative, actual quotient transitions, exact cyclic
pullback/injection, derivative square, unit rule, full-ideal case, and the all-degree attained
monomial ceiling. Existing SplitZero support maps lift their coefficient maps at the original
labels. No duplicate scalar construction or alternative homology core is introduced.

The exact chi\_r annihilator, local basis/multinomial identification and the retraction in Sections
3-5 are supplied here as complete written arguments. Their finite calibrations are separate Python
tests. The new upload's analytic cyclic convolution measure, canonical Gram, differentiated moments,
rank-two control and full trace comparison remain its separately stated inputs and targets. This
work does not independently certify that archive, an analytic moment enclosure, or a tensor-uniform
small-control bound.

\subsubsection{Reproduction}\label{conormal_cyclic_pr21:reproduction}

Run the added \nolinkurl{SplitZero\ conormal\ cyclic\ tower} workflow on this branch. It installs
unchanged Lean 4.31.0 and Mathlib \nolinkurl{fabf563a7c95a166b8d7b6efca11c8b4dc9d911f}, rebuilds all
inherited modules, checks every new source with \nolinkurl{-\/-trust=0\ -DwarningAsError=true}, imports
the combined sources, and rejects any selected transitive axiom outside propext, Classical.choice
and Quot.sound. It reuses the previously tested axiom parser. The original core's byte pin is
checked by the inherited prepare script.

The new finite checker directly multiplies the nilpotent sum in explicit local monomial quotients.
It tests attainment and its literal multinomial coefficient, collided maxima, derivative
divisibility, local dimensions and the h=1 branch, in normal and optimized Python with
deliberate-failure controls. These finite tests do not replace the general Lean proofs or the
written polynomial arguments. The final PR record identifies the exact completed run; no pending run
is a certificate.
